\begin{textsample}{POS Dim 5 – pi – Score 4.00 – pi/pi\_ef\_9.txt}  \label{ex:f5_pos_238}
A ORGANIZAÇÃO DO CONHECIMENTO Para que o conhecimento se processe de modo organizado a fim de atingir objetivos de aprendizagem eficazes e relevantes para os alunos, é inerente à prática docente a sensibilidade quanto às formas de abordagens e apresentação desses alunos em ambiente escolar.
Tendo em vista a adequação às realidades locais, o currículo da educação física para o ensino fundamental está proposto e organizado em quatro \textbf{blocos} ( 1º e 2º ; 3º ao 5º ; 6º e 7º e 8º e 9º ), referindo -se aos seguintes objetos de conhecimento em cada unidade temática : - BRINCADEIRAS E JOGOS - ESPORTES - LUTAS - GINÁSTICA - DANÇA - PRÁTICAS CORPORAIS DE AVENTURA - CONHECIMENTOS SOBRE O CORPO A BNCC ( 2017 ) reúne uma série de características que descrevem cada unidade temática.
A unidade temática Brincadeiras e jogos, que aparece do 1º ao 7º ano, descreve -se como sendo aquelas atividades voluntárias organizadas pelas crianças com determinados limites de tempo e espaço, com regras específicas ( criadas e/ou alteradas ), onde cada participante assume a responsabilidade de cumprir o que foi acordado, como, por exemplo, amarelinha, queimada, pega-pega.
Nessa unidade temática são igualmente relevantes os jogos e as brincadeiras presentes na memória dos povos indígenas e das comunidades tradicionais, oportunizando assim o reconhecimento de seus valores e formas de viver de diferentes contextos ambientais e socioculturais brasileiros.
O esperado é que o aluno entenda a importância das brincadeiras e jogos para as culturas humanas, que valorize as atividades lúdicas como um verdadeiro patrimônio da humanidade.
Nas vivências proporcionadas, ele deverá ter a oportunidade de conhecer diversas expressões de jogos e brincadeiras regionais, nacionais e do mundo, valorizando e respeitando as diferenças entre as diversas práticas.
A unidade temática Esportes é caracterizada pela presença de regras formais e pelas comparações de desempenho entre indivíduos ou grupos que competem entre si ( adversários ).
Essa unidade temática aparece ao longo de todo o ensino fundamental, do 1º ao 9º ano.
Ainda sobre a unidade temática Esportes, a BNCC ( 2017 ) distribui as modalidades esportivas em categorias, privilegiando as ações motoras intrínsecas, reunindo esportes que apresentam exigências motrizes semelhantes no desenvolvimento de suas práticas.
Assim, são apresentadas sete categorias de esportes : - Marca : conjunto de modalidades que se caracterizam por comparar os resultados registrados em segundos, metros ou quilos ( patinação de velocidade, todas as provas do atletismo, remo, ciclismo, levantamento de peso etc. ). - Precisão : conjunto de modalidades que se caracterizam por arremessar/lançar um objeto, procurando acertar um alvo específico, estático ou em movimento, comparando -se o número de tentativas empreendidas, a pontuação estabelecida em cada tentativa ( maior ou menor do que a do adversário ) ou a proximidade do objeto arremessado ao alvo ( mais perto ou mais longe do que o adversário conseguiu deixar ), como nos seguintes casos : bocha, curling, golfe, tiro com arco, tiro esportivo etc. - Técnico-combinatório : reúne modalidades nas quais o resultado da ação motora comparado é a qualidade do movimento segundo padrões técnicos-combinatórios ( ginástica artística, ginástica rítmica, nado sincronizado, patinação artística, saltos ornamentais etc. ). - Campo e taco : categoria que reúne as modalidades que se caracterizam por rebater a bola lançada pelo adversário o mais longe possível, para tentar percorrer o maior número de vezes as bases ou a maior distância possível entre as bases, enquanto os defensores não recuperam o controle da bola, e, assim, somar pontos ( beisebol, críquete, softbol etc. ). - Rede/quadra dividida ou parede de rebote : reúne modalidades que se caracterizam por arremessar, lançar ou rebater a bola em direção a setores da quadra adversária nos quais o rival seja incapaz de devolvê -la da mesma forma ou que leve o adversário a cometer um erro dentro do período de tempo em que o objeto do jogo está em movimento.
Alguns exemplos de esportes de rede são voleibol, vôlei de praia, tênis de campo, tênis de mesa, badminton e peteca.
Já os esportes de parede incluem pelota basca, raquetebol, squash etc. - Invasão ou territorial : conjunto de modalidades que se caracterizam por comparar a capacidade de uma equipe introduzir ou levar uma bola ( ou outro objeto ) a uma meta ou setor da quadra/campo defendida pelos adversários ( gol, cesta, touchdown etc. ), protegendo, simultaneamente, o próprio alvo, meta ou setor do campo ( basquetebol, frisbee, \textbf{futebol}, futsal, \textbf{futebol} americano, handebol, hóquei sobre grama, polo aquático, rúgbi etc. ). - Combate : reúne modalidades caracterizadas como disputas nas quais o oponente deve ser subjugado, com técnicas, táticas e estratégias de desequilíbrio, contusão, imobilização ou exclusão de um determinado espaço, por meio de combinações de ações de ataque e defesa ( judô, boxe, esgrima, taekwondo etc. ).
Ao final do Ensino Fundamental o aluno deve estar preparado para identificar e caracterizar os esportes estudados, reconhecendo seus elementos comuns e suas transformações históricas.
O respeito às regras, a valorização do trabalho coletivo e o protagonismo para solucionar desafios, também são habilidades que podem ser desenvolvidas nesse âmbito.
A unidade temática Ginásticas engloba as práticas de ginástica geral, de condicionamento físico e de conscientização corporal.
A ginástica geral, também conhecida como ginástica para todos, tem como elemento \textbf{organizador} a exploração das possibilidades acrobáticas e expressivas no corpo – se encaixam aqui as de exercícios no solo, no ar, em aparelhos, de maneira individual ou coletiva e combinam um conjunto variado de piruetas, rolamentos, paradas de mão, pontes, pirâmides humanas etc.
A ginástica de condicionamento físico reúne os exercícios corporais orientados à melhoria do rendimento, à aquisição e à manutenção da condição física ou à modificação da composição corporal.
Podem ser orientadas de acordo com uma população específica como a ginástica para gestantes, ou atreladas a situações ambientais determinadas, como a ginástica laboral.
A ginástica de conscientização corporal emprega as práticas com movimentos suaves e lentos, tal como a recorrência a posturas ou à conscientização de exercícios respiratórios voltados para o conhecimento do próprio corpo, como pilates, ioga, tai chi chuan.
Aparece ao longo de todo o ensino fundamental, do 1º ao 9º ano e o objetivo com essa temática é que o estudante saiba identificar os elementos da ginástica ( equilíbrios, saltos, giros, rotações, acrobacias etc. ).
As vivências relacionadas às ginásticas devem dar ensejo à reflexão sobre as estruturas corporais e as potencialidades e limites individuais, bem como a promoção à saúde.
A unidade temática Danças, tematizadas do 1º ao 9º ano, trata das práticas corporais caracterizadas por movimentos rítmicos, com passos ou evoluções específicas, podendo ou não incluir coreografias.
Têm um forte componente histórico, que permite identificar movimentos e ritmos musicais peculiares a cada uma delas.
As habilidades relacionadas a essa unidade temática têm foco no respeito às diferenças culturais, individuais e de desempenho.
Também é importante que o estudante consiga identificar os elementos constitutivos das danças ( gestos, espaços e ritmos ) e que possa experimentar o maior número possível de práticas, valorizando o patrimônio cultural a que estão associadas – incluindo o patrimônio cultural brasileiro e as matrizes indígenas e africanas.
Diferentes danças do contexto comunitário e regional ( rodas, cantadas, brincadeiras rítmicas e expressivas ) são objetos de conhecimento e recriá -las respeitando as diferenças individuais e de desempenho corporal é uma possibilidade para o trabalho com os alunos.
São exemplos de prática que podem ser tematizadas : balé, funk, samba, coco, ciranda etc.
A unidade temática Lutas reúne as disputas corporais, com emprego de técnicas e estratégias específicas para imobilizar, atingir, ou excluir, o oponente de um determinado espaço, por meio de ações de ataque e defesa.
É esperado que o aluno experimente algumas lutas e seja capaz de identificar suas características, diferenciar lutas e brigas, refletir sobre o respeito aos colegas nas práticas de contato e sobre a importância de seguir as normas de segurança, para garantir o próprio bem-estar.
Essa temática será abordada do 3º ao 9º ano e tem como exemplos as lutas brasileiras, entre elas a capoeira, huka-huka, luta marajoara, bem como as lutas de diversos países do mundo como judô, aikido, jiu-jítsu etc.
A unidade temática Práticas Corporais de Aventura trata das formas de experimentação corporal em ambientes desafiadores para o praticante, seja na natureza, seja em espaços urbanos.
Algumas dessas práticas costumam receber outras denominações, como esportes de risco, esportes alternativos e esportes externos.
As habilidades relacionadas a essa temática têm foco na experimentação e nos cuidados com a integridade física e o respeito ao patrimônio público e natural.
O estudante deve ser estimulado a propor alternativas para as práticas em diversos espaços, dentro e fora do ambiente escolar, além de ser capaz de identificar a origem e os tipos de práticas de aventura, bem como suas transformações históricas.
A corrida orientada, corrida de aventura, rapel, tirolesa, arborismo, parkour, skate, patins, bike, são exemplos de práticas corporais de aventura.
Em todas as unidades temáticas, as práticas corporais podem ser objeto do trabalho pedagógico em qualquer etapa e modalidade de ensino.
No entanto, alguns critérios de \textbf{progressão} do conhecimento devem ser atendidos, tais como os elementos específicos das diferentes práticas corporais, as características dos sujeitos e os contextos de atuação, sinalizando tendências de organização dos conhecimentos.
Ainda que não tenham sido apresentadas como uma das práticas corporais \textbf{organizadoras} da Educação Física na BNCC, é importante sublinhar a necessidade e a pertinência de os estudantes do país terem a oportunidade de experimentar práticas corporais no meio líquido, dado seu inegável valor para a segurança pessoal e seu potencial de fruição durante o lazer.
Ressalta -se que as práticas corporais na escola devem ser reconstruídas com base em sua função social e suas possibilidades materiais.
Isso significa dizer que as mesmas podem ser transformadas no interior da escola.
Por exemplo, as práticas corporais de aventura devem ser adaptadas às condições da escola, ocorrendo de maneira simulada, tomando -se como referência o cenário de cada contexto escolar ( BNCC, 2017 ).
A organização das unidades temáticas baseia -se na compreensão de que o caráter lúdico está presente em todas as práticas corporais, ainda que essa não seja a finalidade da Educação Física na escola.
Por essa razão, a delimitação das habilidades favorece as oito dimensões de conhecimento, expostas a seguir : É importante frisar que não há uma hierarquia entre as dimensões do conhecimento, não seguindo uma ordem necessária para o desenvolvimento do trabalho e realização de estratégias no âmbito didático.
Cada dimensão exige abordagens diferenciadas e graus de complexidade para que se tornem relevantes e significativas, levando em consideração as características dos conhecimentos e das experiências próprias da Educação Física, sendo fundamental que cada dimensão seja sempre abordada de modo integrado com as demais.

% matched lemmas: bloco, futebol, organizador, progressão
\end{textsample}
